% Automated Path Search (Ch 18)
% Lines 2761-3617 of main.tex
% This file is for agent editing — will be merged back

\chapter{Automated Path Search}
\label{ch:path-search}

This is the algorithmic heart of CTTT.  Given two Z3 terms $\den{f}$
and $\den{g}$ that are not syntactically equal, the \emph{path search
engine} automatically discovers a multi-step cubical path
$\den{f} \xrightarrow{p_1} h_1 \xrightarrow{p_2} h_2
\xrightarrow{p_3} \ldots \xrightarrow{p_k} \den{g}$
where each $p_i$ is a single path constructor application.

\section{The Path Search Problem}

\begin{definition}[Path Search]
Given Z3 terms $s, t : \PyObj$ (the source and target), the
\emph{path search problem} is:
\[
  \text{Find } p_1, \ldots, p_k \in \mathsf{PathConstructors}
  \text{ such that } s \xrightarrow{p_1} \cdot \xrightarrow{p_2}
  \cdots \xrightarrow{p_k} t
\]
or prove that no such sequence exists (by finding a
counterexample --- a concrete input where $s \neq t$).
\end{definition}

The path constructors are the axioms from the HIT $\mathsf{PyComp}$
(Definition~\ref{def:pycomp}): commutativity, associativity,
fold--unfold, $\beta$-reduction, etc.

\section{The Path Graph}

\begin{definition}[Path Graph]
The \emph{path graph} $G_\mathsf{path}$ has:
\begin{itemize}
  \item \textbf{Vertices}: equivalence classes of Z3 terms modulo
    syntactic identity.
  \item \textbf{Edges}: for each path constructor $p$ and each
    applicable subterm, an edge from the term to the rewritten term.
\end{itemize}
The path search is a graph search from $s$ to $t$ in $G_\mathsf{path}$.
\end{definition}

The path graph is infinite (Z3 terms form an infinite set), but
the search is bounded because:
\begin{enumerate}
  \item Each path constructor application either \emph{simplifies}
    the term (reduces its size/complexity) or \emph{normalizes} it
    (moves it toward a canonical form).
  \item A confluent and terminating rewrite system has a unique
    normal form per equivalence class --- if both $s$ and $t$ reduce
    to the same normal form, a path exists.
\end{enumerate}

\section{The Normalization Strategy}

Rather than searching the full path graph, we normalize both terms
and check if the normal forms coincide:

\begin{definition}[Canonical Normal Form]
\label{def:cnf}
The \emph{canonical normal form} of a Z3 term $t$ is obtained by
exhaustively applying the following rewrite rules (in priority order):

\textbf{Phase 1: $\beta$-reduction (definitional equality).}
\begin{enumerate}
  \item Inline all nested function applications
    ($(\lambda x.\, B)(a) \to B[x := a]$).
  \item Unfold all \texttt{\_\_func\_\_} references.
  \item Collapse trivial \texttt{If} branches
    ($\mathsf{If}(\top, a, b) \to a$; $\mathsf{If}(\bot, a, b) \to b$).
\end{enumerate}

\textbf{Phase 2: Arithmetic normalization (ring axioms).}
\begin{enumerate}
  \setcounter{enumi}{3}
  \item Commute operands to canonical order
    ($b + a \to a + b$ when $a <_\mathsf{lex} b$).
  \item Associate right ($((a+b)+c) \to (a+(b+c))$).
  \item Eliminate identities ($a + 0 \to a$; $a \times 1 \to a$).
  \item Absorb zeros ($a \times 0 \to 0$).
  \item Cancel double negation ($-(-a) \to a$).
\end{enumerate}

\textbf{Phase 3: Fold canonicalization (Nat-recursor).}
\begin{enumerate}
  \setcounter{enumi}{8}
  \item Detect recursive descent patterns and replace with
    canonical $\fold(\op, \mathsf{start}, \mathsf{stop}, \mathsf{acc})$.
  \item Detect accumulation loops and replace with canonical $\fold$.
  \item When both rec and iter produce $\fold$ terms with the same
    $(\op, \mathsf{start}, \mathsf{stop}, \mathsf{acc})$, they
    are syntactically identical.
\end{enumerate}

\textbf{Phase 4: Higher-order normalization (operadic).}
\begin{enumerate}
  \setcounter{enumi}{11}
  \item Apply map--fold fusion
    ($\mathsf{reduce}(\oplus, e, \mathsf{map}(h, xs)) \to
    \mathsf{reduce}(\lambda a,x.\, a \oplus h(x), e, xs)$).
  \item Apply comprehension canonicalization
    ($[e(x) \text{ for } x \text{ in } xs] \to
    \mathsf{comp}_{\mathsf{hash}(e)}(xs)$).
  \item Apply idempotence ($\mathsf{sorted}(\mathsf{sorted}(x)) \to
    \mathsf{sorted}(x)$).
  \item Apply involutions ($\mathsf{reversed}(\mathsf{reversed}(x))
    \to x$).
\end{enumerate}

\textbf{Phase 5: Shared symbol unification (univalence).}
\begin{enumerate}
  \setcounter{enumi}{15}
  \item Replace all builtin calls with shared function symbols
    ($\mathsf{sorted}_k(x) \to \mathsf{py\_sorted}(x)$).
  \item Replace all method calls with shared method symbols
    ($\mathsf{meth\_append}_k(o, v) \to
    \mathsf{py\_meth\_append}(o, v)$).
\end{enumerate}
\end{definition}

\begin{theorem}[Normalization Confluence]
\label{thm:confluence}
The rewrite system of Definition~\ref{def:cnf} is \emph{locally
confluent}: if $t \to t_1$ and $t \to t_2$ by different rules,
there exists $t'$ such that $t_1 \to^* t'$ and $t_2 \to^* t'$.

Combined with termination (each rule decreases a complexity measure),
the system is \emph{confluent} (by Newman's lemma): every term has
a unique normal form.
\end{theorem}

\begin{proof}[Proof sketch]
The phases are ordered by priority, and rules within each phase are
orthogonal (they apply to disjoint subterms).  Critical pairs arise
only between commutativity and associativity, which are resolved by
the canonical ordering $<_\mathsf{lex}$.
\end{proof}

\begin{corollary}[Path Existence from Normal Forms]
\label{cor:path-nf}
A cubical path from $\den{f}$ to $\den{g}$ exists if and only if
their canonical normal forms are identical:
\[
  \mathsf{nf}(\den{f}) = \mathsf{nf}(\den{g})
  \iff
  \exists p : \Path(\den{f}, \den{g})
\]
The path is reconstructed by composing the rewrite steps:
$\den{f} \to^* \mathsf{nf}(\den{f}) = \mathsf{nf}(\den{g})
\leftarrow^* \den{g}$.
\end{corollary}

\section{The Path Search Algorithm}

\begin{lstlisting}[caption={Automated path search}]
def search_path(term_f, term_g, theory):
    """Search for a cubical path from term_f to term_g.

    Returns:
      PathResult(found=True, path=[step1, step2, ...])
      or PathResult(found=False, reason=...)
    """
    # Phase 1-5: normalize both terms
    nf_f, steps_f = normalize(term_f, theory)
    nf_g, steps_g = normalize(term_g, theory)

    # Check syntactic equality of normal forms
    if z3_syntactic_eq(nf_f, nf_g):
        # Path found! Compose forward and backward steps.
        path = steps_f + reversed(steps_g)
        return PathResult(found=True, path=path)

    # Normal forms differ: try Z3 semantic check
    # (maybe the terms are equal but our rewrite system
    # doesn't know the right axioms)
    solver = Solver()
    theory.semantic_axioms(solver)
    solver.add(nf_f != nf_g)
    result = solver.check()

    if result == unsat:
        # Z3 proved equality using axioms we didn't rewrite
        return PathResult(found=True,
                          path=steps_f + ['z3_axiom'] + reversed(steps_g))
    elif result == sat:
        # Genuine counterexample: no path exists
        return PathResult(found=False,
                          counterexample=solver.model())
    else:
        return PathResult(found=None, reason='timeout')
\end{lstlisting}

\section{Path Search on Each Fiber}

The path search is performed \emph{per fiber} (per type assignment),
not globally.  This is the cubical--cohomological interaction:

\begin{lstlisting}[caption={Fiber-wise path search}]
def fiber_wise_path_search(secs_f, secs_g, sites, theory):
    """Search for paths on each fiber independently.

    Returns per-fiber path results and the global H^1.
    """
    local_paths = {}

    for combo, site_id in sites.items():
        # Restrict terms to this fiber
        fiber_f = restrict_to_fiber(secs_f, combo, theory)
        fiber_g = restrict_to_fiber(secs_g, combo, theory)

        # Search for path on this fiber
        for sf in fiber_f:
            for sg in fiber_g:
                if not guards_overlap(sf.guard, sg.guard):
                    continue
                result = search_path(sf.term, sg.term, theory)
                if result.found == False:
                    return GlobalResult(equiv=False,
                        fiber=combo, cx=result.counterexample)
                elif result.found == True:
                    local_paths[site_id] = result.path

    # All fibers have paths: check gluing via Cech H^1
    h1 = compute_cech_h1(local_paths, sites)
    if h1 == 0:
        return GlobalResult(equiv=True, paths=local_paths)
    else:
        return GlobalResult(equiv=False, obstruction=h1)
\end{lstlisting}

\section{Path Composition and the Proof Certificate}

When the path search succeeds, it produces a \emph{proof certificate}:
a sequence of named path constructor applications that transforms
$\den{f}$ into $\den{g}$.

\begin{definition}[Proof Certificate]
A \emph{proof certificate} for $f \equiv g$ is a list of triples:
\[
  [(r_1, \mathsf{pos}_1, \mathsf{ax}_1),\;
   (r_2, \mathsf{pos}_2, \mathsf{ax}_2),\;
   \ldots,\;
   (r_k, \mathsf{pos}_k, \mathsf{ax}_k)]
\]
where $r_i$ is the term at step $i$, $\mathsf{pos}_i$ is the subterm
position, and $\mathsf{ax}_i$ is the axiom applied.  The certificate
satisfies:
\begin{align}
  r_0 &= \den{f} \\
  r_{i+1} &= r_i[\mathsf{pos}_i := \mathsf{ax}_i(r_i|_{\mathsf{pos}_i})] \\
  r_k &= \den{g}
\end{align}
\end{definition}

\begin{example}[Certificate for Factorial Equivalence]
\begin{align}
  \den{f} &= \mathsf{rec\_f}(p_0)
    && \text{(recursive definition)} \\
  \xrightarrow{\mathsf{D1},\, \text{root}} &\;
    \fold(\mathsf{MUL}, 1, \mathsf{ival}(p_0)+1, \mathsf{IntObj}(1))
    && \text{(Nat-eliminator extraction)} \\
  &= \den{g}
    && \text{(same as iterative fold)}
\end{align}
Certificate: $[(\mathsf{rec\_f}(p_0), \epsilon, \mathsf{D1})]$.
One step, at the root, using the Nat-recursor path.
\end{example}

\begin{example}[Multi-Step Certificate]
Programs that compute $\mathsf{sum}(\mathsf{reversed}(\mathsf{sorted}(x)))$
vs.\ $\mathsf{sum}(x)$:
\begin{align}
  &\mathsf{py\_sum}(\mathsf{py\_reversed}(\mathsf{py\_sorted}(p_0))) \\
  \xrightarrow{\mathsf{reversed\_invol},\, [1,1]} &\;
    \mathsf{py\_sum}(\mathsf{py\_sorted}(p_0))
    && \text{(sorted absorbs reversed)} \\
  \xrightarrow{\mathsf{sum\_perm},\, [1]} &\;
    \mathsf{py\_sum}(p_0)
    && \text{(sum is perm-invariant)}
\end{align}
Certificate: $[(\ldots, [1,1], \mathsf{sorted\_absorbs\_reversed}),\;
(\ldots, [1], \mathsf{sum\_perm\_invariant})]$.
\end{example}

\section{Path Search Strategies}

The basic normalization strategy (Definition~\ref{def:cnf}) is
``normalize both sides and compare.''  But more sophisticated
strategies exist:

\subsection{Bidirectional Search}

Search from both $\den{f}$ and $\den{g}$ simultaneously,
meeting in the middle:

\begin{lstlisting}[caption={Bidirectional path search}]
def bidirectional_search(s, t, axioms, max_depth=5):
    forward = {hash(s): (s, [])}    # term -> (term, path)
    backward = {hash(t): (t, [])}

    for depth in range(max_depth):
        # Expand forward frontier
        new_forward = {}
        for h, (term, path) in forward.items():
            for ax in axioms:
                for pos in applicable_positions(term, ax):
                    new_term = apply_axiom(term, ax, pos)
                    nh = hash(new_term)
                    if nh in backward:
                        # FOUND: path from s to t
                        _, bpath = backward[nh]
                        return path + [(term, pos, ax)] + \
                               reversed(bpath)
                    new_forward[nh] = (new_term,
                        path + [(term, pos, ax)])
        forward.update(new_forward)

        # Similarly expand backward frontier
        # ...

    return None  # no path within depth bound
\end{lstlisting}

\subsection{Fiber-Guided Search}

Use the \v{C}ech decomposition to \emph{guide} the search:
\begin{enumerate}
  \item Identify which fibers are ``easy'' (terms already match or
    nearly match after normalization).
  \item On hard fibers, use the easy fibers' solutions as hints:
    if the path on the Int fiber used axiom $\mathsf{D1}$, try
    $\mathsf{D1}$ on the Str fiber first.
  \item This is \emph{transfer along the fibration}: a successful
    path on one fiber is transported to adjacent fibers.
\end{enumerate}

\subsection{Axiom Relevance Filtering}

Not all axioms are relevant for a given pair.  Filter by:
\begin{enumerate}
  \item \textbf{Symbol filtering}: only use axioms mentioning
    symbols that appear in $\den{f}$ or $\den{g}$.
  \item \textbf{Sort filtering}: only use axioms whose type
    signature matches the current fiber.
  \item \textbf{Depth filtering}: prioritize axioms that reduce
    term depth (simplifying rules before complexifying ones).
\end{enumerate}

\section{Sheaf Descent for Tractability of Undecidable Problems}
\label{sec:descent-tractability}

Program equivalence is undecidable in general (Rice's theorem).
Specification checking and bug finding reduce to it, so they are
also undecidable.  Yet we claim CTTT makes these problems
\emph{tractable} for a large class of real programs.  This section
explains \emph{why} --- the mathematical mechanism by which sheaf
descent converts undecidable global problems into decidable local ones.

\subsection{The Undecidability Barrier}

The global equivalence query is:
\[
  \forall x : \PyObj.\; \den{f}(x) = \den{g}(x)
\]
This is a universal quantification over an infinite domain ($\PyObj$
has infinitely many inhabitants via $\mathsf{IntObj}(n)$ for all
$n \in \Int$).  No algorithm can decide this in general.

But the undecidability proof (Rice's theorem) requires adversarial
programs --- programs specifically constructed to defeat any analysis.
Real programs are not adversarial.  They have \emph{structure} that
can be exploited.

\subsection{How Descent Exploits Structure}

Sheaf descent exploits two kinds of structure:

\textbf{1. Type structure (the fibration).}
Real Python programs dispatch on types --- they use
\texttt{isinstance}, duck typing, different code paths for different
input types.  This means the universal quantification decomposes
into a \emph{finite} disjunction:
\[
  \forall x.\; \den{f}(x) = \den{g}(x)
  \iff
  \bigwedge_{\tau \in \mathsf{TypeTag}} \forall x.\;
    \mathsf{tag}(x) = \tau \implies \den{f}(x) = \den{g}(x)
\]

Each conjunct is a \emph{fiber-restricted} query.  The finite
conjunction is decidable if each conjunct is.

\textbf{2. Algebraic structure (the path constructors).}
Within each fiber, the Z3 terms have algebraic structure:
commutativity, associativity, fold patterns.  The path constructors
(equational axioms) transform this structure, often collapsing
complex terms to simple canonical forms.

The canonical form is unique (by confluence, Theorem~\ref{thm:confluence}),
and comparing canonical forms is decidable (syntactic equality check).

\subsection{The Descent Decomposition Theorem}

\begin{theorem}[Descent Makes Equivalence Tractable]
\label{thm:descent-tractable}
Let $\covers = \{U_\tau\}_{\tau \in T}$ be the covering of
$\PyObj$ by type fibers, where $T$ is a finite set of type tags.
Then:
\begin{enumerate}
  \item The global equivalence query
    $\forall x.\, \den{f}(x) = \den{g}(x)$ reduces to
    $|T|$ local queries plus one $\Hone$ computation.
  \item Each local query is in the theory QF\_DT + QF\_LIA + QF\_S
    (quantifier-free datatypes + linear integer arithmetic + strings),
    which is \emph{decidable}.
  \item The $\Hone$ computation is decidable ($O(|T|^3)$ via GF(2)
    rank).
  \item Therefore: the full equivalence check is decidable for
    programs whose Z3 denotations fall in the quantifier-free
    fragment (no RecFunctions, no universally quantified axioms).
\end{enumerate}
\end{theorem}

\begin{proof}
(1) By the descent theorem (\ref{thm:descent-full}), the global query
decomposes into local queries on each fiber plus the gluing check
($\Hone = 0$).

(2) On each fiber $U_\tau$, the type is fixed (e.g., $\mathsf{tag}(x) =
\mathsf{TInt}$ forces $x = \mathsf{IntObj}(n)$ for some $n$).
The Z3 terms reduce to expressions over $\mathsf{IntSort}$
(decidable QF\_LIA), $\mathsf{BoolSort}$ (propositional, decidable),
$\mathsf{StringSort}$ (decidable QF\_S), or $\mathsf{PyObj}$ constructors
(decidable QF\_DT).  The combined theory QF\_DT + QF\_LIA + QF\_S
is decidable by Nelson--Oppen combination.

(3) GF(2) linear algebra is decidable, $O(n^3)$.

(4) Combining (1)--(3): decidable $+$ decidable $+$ decidable $=$
decidable.
\end{proof}

\begin{remark}[Beyond the Quantifier-Free Fragment]
When programs use RecFunctions (from loops/recursion), the local
queries may exit the QF fragment.  The Nat-eliminator extraction
(Chapter~\ref{ch:nat-elim}) recovers decidability for many such
programs by converting RecFunctions to canonical folds, which are
ground terms (no quantifiers).

When equational axioms involve universal quantifiers ($\forall x.\,
\mathsf{sorted}(\mathsf{sorted}(x)) = \mathsf{sorted}(x)$), the
local queries are semi-decidable (Z3 may timeout).  But the axioms
are applied \emph{as rewrite rules} during normalization, not as
quantified constraints --- so they terminate (the rewrite system is
terminating by Theorem~\ref{thm:confluence}).
\end{remark}

\subsection{Complexity Analysis}

\begin{theorem}[Complexity of Descent-Based Equivalence]
\label{thm:descent-complexity}
For programs with $n$ parameters, each having at most $k$ possible
type tags, the descent-based equivalence check has complexity:
\[
  O(k^n \cdot C_\mathsf{local} + k^{3n})
\]
where $C_\mathsf{local}$ is the cost of a single local Z3 query.

For the common case of single-typed parameters ($k = 1$):
\[
  O(C_\mathsf{local} + 1) = O(C_\mathsf{local})
\]
--- a single Z3 query, no combinatorial blowup.

For the worst case ($k = 7$ types, $n$ parameters):
$k^n = 7^n$ fibers --- exponential in the number of parameters.
But duck type inference typically reduces $k$ to 1--2 per parameter,
so the practical complexity is $O(2^n \cdot C_\mathsf{local})$.
\end{theorem}

\subsection{Comparison with Prior Approaches}

\begin{center}
\small
\begin{tabular}{l l l l}
  \textbf{Approach} & \textbf{Global query} & \textbf{Decidable?} &
  \textbf{Practical?} \\
  \hline
  Direct Z3 & $\forall x.\, f(x) = g(x)$ & Semi-decidable &
    Slow/timeout \\
  Testing & Sample $x_1, \ldots, x_N$ & Incomplete &
    Fast but unsound \\
  Abs.\ interp. & $\alpha(f) = \alpha(g)$ & Decidable &
    Fast but imprecise \\
  CTTT descent & $\bigwedge_\tau (\forall x_\tau.\, f = g) \land
    \Hone = 0$ & \textbf{Decidable (QF)} & \textbf{Fast \& sound} \\
\end{tabular}
\end{center}

CTTT descent is the only approach that is simultaneously decidable
(for the QF fragment), sound (no false positives/negatives within
the modeled domain), and practical (scales with duck-type narrowing).

\subsection{The Refinement Principle}

When descent alone is insufficient (the local Z3 query times out or
returns ``unknown''), we can \emph{refine} the covering:

\begin{definition}[Covering Refinement]
A refinement $\covers' = \{V_j\}$ of covering $\covers = \{U_i\}$
is a covering such that each $V_j \subseteq U_i$ for some $i$.
Refinement makes fibers \emph{smaller} (more constrained), which
makes local queries \emph{easier}.
\end{definition}

\begin{theorem}[Refinement Soundness]
\label{thm:refinement}
If the equivalence check succeeds on a refinement $\covers'$
($\Hone(\covers', \Iso) = 0$), it also succeeds on the original
covering ($\Hone(\covers, \Iso) = 0$).

The converse is NOT true: a refinement may introduce new overlaps
and new transition checks, which could fail even if the original
covering succeeds.
\end{theorem}

\begin{proof}
The refinement morphism $\covers' \to \covers$ induces a morphism
of \v{C}ech complexes $C^\bullet(\covers) \to C^\bullet(\covers')$.
By functoriality of cohomology, $\Hone(\covers) = 0$ implies
$\Hone(\covers') = 0$ (the image of a trivial class is trivial).
\end{proof}

Refinement strategies:
\begin{enumerate}
  \item \textbf{Value-range refinement}: split the Int fiber into
    $\{n \geq 0\}$ and $\{n < 0\}$ (positive and negative integers).
    This helps with programs that have different behavior for positive
    vs.\ negative inputs.
  \item \textbf{Size refinement}: split the Pair fiber into
    $\{|xs| = 0\}$, $\{|xs| = 1\}$, $\{|xs| \geq 2\}$ (empty,
    singleton, longer lists).  This helps with programs that have
    base cases for short inputs.
  \item \textbf{Predicate refinement}: split a fiber by a predicate
    derived from the program's branch conditions.  E.g., if $f$
    branches on $n \leq 1$, split the Int fiber at $n = 1$.
\end{enumerate}

\begin{theorem}[Predicate Refinement from Branch Conditions]
\label{thm:pred-refinement}
Let $\mathsf{guards}_f = \{g_1, \ldots, g_m\}$ and
$\mathsf{guards}_g = \{h_1, \ldots, h_p\}$ be the branch conditions
of programs $f$ and $g$ respectively.  The \emph{branch-condition
refinement} is the covering:
\[
  \covers_\mathsf{branch} = \{U \cap R \mid U \in \covers,\;
    R \in \mathsf{atoms}(\{g_i\} \cup \{h_j\})\}
\]
where $\mathsf{atoms}$ are the minimal non-empty regions of the
partition induced by the guards.

On each atom, \emph{both} programs follow a single code path
(no branching), so their Z3 terms are ground expressions
(no $\mathsf{If}$), which are decidable.
\end{theorem}

\begin{proof}
Within an atom $R$, all guard conditions have fixed truth values.
Therefore $\mathsf{If}(\mathsf{guard}, t, f)$ simplifies to $t$ or
$f$ (no conditional).  The resulting Z3 terms are quantifier-free
ground expressions over the base theories, which are decidable.
\end{proof}

\begin{corollary}[Branch-Refined Descent is Decidable]
The branch-condition refinement produces decidable local queries
for \emph{any} program (not just QF programs).  The only remaining
source of undecidability is the RecFunction bodies (loops/recursion).
\end{corollary}

This is a powerful result: by refining the covering using the
programs' own branch conditions, we eliminate branching complexity
from the local queries, leaving only loop/recursion complexity.

\subsection{Iterated Descent: The Descent Tower}

For programs with nested structure (loops inside branches inside
loops), we can apply descent \emph{iteratively}:

\begin{definition}[Descent Tower]
\label{def:descent-tower}
A \emph{descent tower} is a sequence of coverings:
\[
  \covers_0 \leftarrow \covers_1 \leftarrow \covers_2
    \leftarrow \ldots \leftarrow \covers_k
\]
where each $\covers_{i+1}$ is a refinement of $\covers_i$.

Level 0: type fibers (Int, Str, List, ...).
Level 1: type fibers $\times$ branch conditions.
Level 2: type fibers $\times$ branch conditions $\times$ loop bounds.
Level 3: type fibers $\times$ branch conditions $\times$ loop bounds
  $\times$ recursion depth.
\end{definition}

At each level, the local queries become simpler (more constrained,
fewer degrees of freedom).  Eventually, the local queries are ground
and decidable.

\begin{theorem}[Descent Tower Convergence]
\label{thm:tower-convergence}
For programs with finitely many branches and bounded-depth
recursion, the descent tower converges at a finite level $k$:
all local queries at level $k$ are ground and decidable.
\end{theorem}

\begin{proof}
Each level refines by a finite partition (finitely many branches,
finitely many loop bound values).  The atoms at level $k$ correspond
to complete execution traces of bounded depth.  On each trace, the
Z3 term is ground (no $\mathsf{If}$, no RecFunction unfolding needed
beyond the bound).
\end{proof}

\begin{remark}[Unbounded Recursion]
For programs with unbounded recursion, the descent tower does NOT
converge.  This is the residual source of undecidability in CTTT.
The Nat-eliminator extraction (Chapter~\ref{ch:nat-elim}) handles
many unbounded cases by recognizing them as folds, but general
recursion remains semi-decidable.
\end{remark}

\subsection{The Effective Descent Condition}

Not every covering admits effective descent.  We characterize when
descent is most useful:

\begin{definition}[Effective Descent]
A covering $\covers$ admits \emph{effective descent} for the
equivalence check $f \equiv g$ if:
\begin{enumerate}
  \item Each local query $\Iso(U_i)$ is decidable (terminates with
    a definite answer).
  \item The number of overlaps is polynomial in $|\covers|$ (the
    \v{C}ech complex is sparse).
  \item The transition checks $\varphi_{ij}$ are all trivial
    (both local paths agree automatically on overlaps).
\end{enumerate}
\end{definition}

\begin{theorem}[When Descent is Effective]
\label{thm:effective-descent}
Descent is effective when:
\begin{enumerate}
  \item \textbf{Duck type inference is sharp}: each parameter has
    a unique type (no polymorphism).  Then $|\covers| = 1$, no
    overlaps, no transitions.  This is the common case.
  \item \textbf{Fibers are independent}: the program's behavior on
    fiber $U_i$ does not depend on behavior on fiber $U_j$.  Then
    transitions are trivial.
  \item \textbf{Branch conditions align}: both programs branch on
    the same conditions.  Then the branch-condition refinement
    produces the same atoms for both.
\end{enumerate}
\end{theorem}

\subsection{Extensions: Beyond Type Fibers}

The descent framework is not limited to type-based coverings.  Any
\emph{partition} of the input space that respects program structure
can serve as a covering:

\begin{enumerate}
  \item \textbf{Value-range covering}: partition by value ranges
    ($n < 0$, $n = 0$, $n > 0$, $n > 100$, etc.).  Useful for
    programs with thresholds.
  \item \textbf{Shape covering}: partition by data structure shape
    (empty list, singleton, doubleton, general list).  Useful for
    recursive programs with base cases.
  \item \textbf{Reachability covering}: partition by reachable
    code paths (which branches are taken).  This is the
    branch-condition refinement from Theorem~\ref{thm:pred-refinement}.
  \item \textbf{Symmetry covering}: partition by symmetry group
    orbits (inputs related by permutation).  Useful for programs
    whose output is permutation-invariant.
  \item \textbf{Norm covering}: partition by norm/size ($|x| \leq k$
    for various $k$).  The spectral sequence
    (Chapter~\ref{ch:spectral}) is the cohomology of this filtration.
\end{enumerate}

Each covering type provides a different decomposition strategy,
and the optimal covering depends on the programs being compared.
The \emph{covering selection heuristic} is:
\begin{enumerate}
  \item Start with type fibers (cheapest, most effective for
    well-typed programs).
  \item If local queries fail, refine by branch conditions.
  \item If still failing, refine by value ranges or shapes.
  \item If still failing, try symmetry or norm coverings.
\end{enumerate}

This hierarchical refinement is the \emph{spectral sequence in
action}: each refinement level corresponds to a page of the
spectral sequence, and the sequence converges when the local
queries become decidable.


\section{Soundness and Completeness}

\begin{theorem}[Path Search Soundness]
If the path search returns a certificate $[r_0, \ldots, r_k]$, then
$r_0 = \den{f}$, $r_k = \den{g}$, and each step $r_i \to r_{i+1}$
is justified by a valid path constructor.  Therefore
$\den{f} \equiv \den{g}$ (the programs are semantically equivalent).
\end{theorem}

\begin{proof}
Each path constructor is a theorem of the cubical type theory
(proved in the relevant dichotomy chapter).  Kan composition
(Chapter~\ref{ch:kan}) ensures the steps compose to a valid path.
\end{proof}

\begin{theorem}[Path Search Completeness (Relative)]
\label{thm:path-completeness}
The path search is \emph{complete relative to the axiom set}: if
a path exists using only the axioms in the HIT $\mathsf{PyComp}$,
the normalization procedure finds it (assuming the rewrite system
is confluent and terminating for the applicable axiom subset).

The path search is \emph{incomplete} in general: there exist
equivalent programs for which no path exists in $\mathsf{PyComp}$
(because the axiom set doesn't cover all valid equalities ---
Rice's theorem prevents completeness).
\end{theorem}

\section{Path Search for Spec Checking}

The same path search engine handles spec checking by searching for
a path from $\den{f}$ to $\den{S}$ (the specification's canonical
form):

\begin{lstlisting}[caption={Spec check via path search}]
def spec_check(program_source, spec_source, theory):
    """Check if program satisfies spec via path search.

    spec_source is a reference implementation or a Z3 formula.
    """
    term_f = compile(program_source, theory)
    term_spec = compile(spec_source, theory)

    result = fiber_wise_path_search(
        sections_of(term_f),
        sections_of(term_spec),
        build_sites(program_source, spec_source, theory),
        theory
    )

    if result.equiv:
        return SpecResult(satisfies=True,
                          certificate=result.paths)
    else:
        return SpecResult(satisfies=False,
                          bug_fiber=result.fiber,
                          counterexample=result.cx)
\end{lstlisting}

\section{Path Search for Bug Finding}

Bug finding is path search with \emph{failure analysis}: when no
path exists on a fiber, the failure point reveals the bug.

\begin{lstlisting}[caption={Bug finding via failed path search}]
def find_bugs(program_source, spec_source, theory):
    """Find bugs by attempting path search on each fiber.

    Returns list of (fiber, depth, counterexample) triples.
    """
    bugs = []
    term_f = compile(program_source, theory)
    term_spec = compile(spec_source, theory)
    sites = build_sites(program_source, spec_source, theory)

    for combo, site_id in sites.items():
        result = search_path(
            restrict(term_f, combo),
            restrict(term_spec, combo),
            theory
        )
        if result.found == False:
            bugs.append(Bug(
                fiber=combo,
                counterexample=result.counterexample,
                # How far did the path get before failing?
                partial_path=result.partial_steps,
                # What axiom would be needed?
                missing_axiom=result.stuck_on
            ))

    return BugReport(bugs=bugs,
                     severity=len(bugs) / len(sites),
                     h1_rank=compute_h1_rank(bugs, sites))
\end{lstlisting}

The \texttt{partial\_path} and \texttt{missing\_axiom} fields provide
\emph{diagnostic information}: they tell the developer not just
\emph{that} the program is wrong, but \emph{why} the proof failed
and \emph{what kind of fix} is needed.

\section{Path-Based Automatic Repair Suggestion}

When the path search fails at a specific step, the failure point
suggests a repair:

\begin{theorem}[Repair from Path Failure]
\label{thm:repair-from-path}
If the path search from $\den{f}$ toward $\den{S}$ succeeds for
$k$ steps and fails at step $k+1$ because axiom $\mathsf{ax}$ is
not applicable:
\[
  \den{f} \xrightarrow{p_1} \cdots \xrightarrow{p_k} h_k
  \xrightarrow{\mathsf{ax}?} \text{STUCK}
\]
then the \emph{repair} is: modify $f$ so that $h_k$ satisfies the
precondition of $\mathsf{ax}$.  The modification is localized to
the subterm at position $\mathsf{pos}_{k+1}$ of $h_k$.
\end{theorem}

\begin{example}
A program tries to compute the sum of a list but uses multiplication
instead of addition:
\begin{lstlisting}
def f(xs):
    acc = 0
    for x in xs: acc *= x  # BUG: should be +=
    return acc
\end{lstlisting}

Path search toward spec $S = \mathsf{py\_sum}(xs)$:
\begin{enumerate}
  \item $\den{f} = \fold(\mathsf{MUL}, \ldots)$ (detected from loop).
  \item Normalize: $\fold(\mathsf{MUL}, 0, \ldots) = 0$ (because
    $0 \times \text{anything} = 0$ by absorbing zero axiom).
  \item Target: $\mathsf{py\_sum}(xs)$ normalizes to
    $\fold(\mathsf{ADD}, 0, \ldots)$.
  \item STUCK: $\fold(\mathsf{MUL}, \ldots) \neq \fold(\mathsf{ADD}, \ldots)$.
    The path constructor D1 applies but the operation tag differs:
    $\mathsf{MUL} \neq \mathsf{ADD}$.
\end{enumerate}

Repair suggestion: ``Change $\mathsf{MUL}$ to $\mathsf{ADD}$ at
the loop accumulation step (line 3: \texttt{*=} $\to$ \texttt{+=}).''
\end{example}


% ══════════════════════════════════════════════════════════
\part{The Twenty-Four Dichotomies}
% ══════════════════════════════════════════════════════════

